\section{Finite prime heads, compact radical repair, and two essential tail signs}
\label{vfourA:sec:prime-obstruction}

This appendix integrates the finite-head obstruction and known-radical repair
of~\cite{vfourPrimeHead} with the positive-range tail obstruction
of~\cite{vfourPositiveRange}.  The first obstruction occurs on actual
zero-energy vectors of the complete form; the second occurs on positive
spectral ranges of the complete operator.  They are distinct statements,
and neither supplies a sign for the complete form on its unresolved sector.

\subsection{Common normalization and inherited analytic inputs}
\label{vfourA:subsec:inputs}

We use the kernel, spaces, and complete forms of
Section~\ref{sec:foundations}.  Inner products in this appendix are linear
in the first argument and conjugate-linear in the second; in particular,
$\ip uv_{\Hh}=\int u\overline v/W$.  Recall
\[
 \begin{gathered}
 c(x)=\cosh(x/2),\qquad W=\frac{\kappa}{2c},\qquad
 \Hh=L^2_{\mathrm{even}}(dx/W),\\
 m(u)=\int_\R c(x)u(x)\,dx,\qquad \Hh_c=\ker m.
 \end{gathered}
\]
The normalization $\int2c\kappa=1$ gives
$\norm\kappa_{\Hh}=1$, $\ip u\kappa_{\Hh}=2m(u)$, and the orthogonal
projection $Cu=u-2m(u)\kappa$ onto $\Hh_c$.  Write
$X=\{u\in\Hh:\tg[u]<\infty\}$ and $X_c=X\cap\Hh_c$.
For an integer $M\ge1$, let
\begin{align}
 p_{\le M}[u]&=2\operatorname{Re}\sum_{2\le n\le M}
 \frac{\Lambda(n)}{\sqrt n}
 \int_\R u(x+\log n)\overline{u(x)}\,dx,\label{vfourA:eq:head}\\
 p_{>M}&=p-p_{\le M},\qquad
 q_M=\tg-\cg\norm{\cdot}_2^2-p_{\le M}
     =q_0+p_{>M}.\label{vfourA:eq:head-complete}
\end{align}
Thus the subscript $0$ continues to denote the \emph{complete} form, not
the empty head.  All prime powers, including all powers below a finite
cutoff, and both shift orientations are retained.  We use the same symbols
for the polarized Hermitian forms.

For clarity, denote the weighted representative $B$ of
Section~\ref{sec:weighted-physical} by $L$ in this appendix:
\begin{equation}\label{vfourA:eq:actual-operator}
 L=T_\Gamma-D_b,\qquad d:=\norm{D_b}<\infty,\qquad
 \dom L=\dom T_\Gamma,\qquad L_c=L|_{\Hh_c}.
\end{equation}
Here $D_b$ is exactly the bounded complete correction in
\eqref{eq:Db-physical}, including every prime power.  The weak source
identity of Proposition~\ref{a:prop:thetanull} gives
$L\kappa=-\kappa/2$, so $C$ reduces $L$.  Lemma~\ref{d:lem:Gamma-H1}
admits even compact smooth functions to $\dom L$ and identifies
$T_\Gamma u=WD_\Gamma u$ there.  The ordinary Fourier multiplier is
\begin{equation}\label{vfourA:eq:symbol}
 d_\Gamma(t)=2\int_0^\infty\rho(s)(1-\cos ts)\,ds
 =\operatorname{Re}\psi(1/4+it/2)-\psi(1/4).
\end{equation}
In particular $d_\Gamma\ge0$ and
$d_\Gamma(t)=\log|t|+O(1)$ at infinity; these follow from the complete
Gamma integral and the digamma identities in~\cite[Sections 5.9 and 5.11]{DLMF}.

The strip continuation and bounds used below are
Lemma~\ref{d:lem:theta-tails}; the exact identity
$\mathcal F\kappa=\Xi=\xi(1/2+i\,\cdot)$ is
\eqref{d:eq:Xi-xi}, with
$\mathcal Fu(t)=\int u(x)e^{itx}\,dx$.  The full-domain radical input is
Lemma~\ref{a:lem:radical}.  In its notation set
\begin{equation}\label{vfourA:eq:radical-spaces}
 r_j=\kappa^{(2j)}-4^{-j}\kappa,\qquad
 \Rcal=\operatorname{span}\{r_j:j\ge1\},\qquad
 \Nn=\overline{\Rcal}^{\Hh},\qquad \Kk=\Hh_c\ominus\Nn.
\end{equation}
Since $q_0(r_j,z)=0$ for every $z\in X$, representation and closedness
give $\Nn\subseteq\ker L_c$.  No equality with the full physical kernel
is asserted.  The ordinary-$L^2$ minimal form realization and compact
resolvent used in the spectral consequence below are precisely
Propositions~\ref{a:prop:closedpencil} and~\ref{a:prop:compact} in
Section~\ref{a:sec:pencil}; they are not statements of unitary equivalence
with $L_c$.

\subsection{Infinite positive and negative index on the actual radical}
\label{vfourA:subsec:index}

Put $w(t)=1+\log(2+|t|)$ and $R(t)=(t^2+1/4)\Xi(t)$.  For a polynomial
$P$, the actual differential radical
\begin{equation}\label{vfourA:eq:polynomial-radical}
 r_P=(-\partial_x^2+1/4)P(-\partial_x^2)\kappa\in\Rcal,
 \qquad \mathcal F r_P(t)=R(t)P(t^2)
\end{equation}
belongs to $\Rcal$ because its differential polynomial in
$-\partial_x^2$ vanishes at $-1/4$.

\begin{lemma}[Weighted Fourier density]\label{vfourA:lem:density}
The functions $R(t)P(t^2)$ are dense in $L^2_{\mathrm{even}}(w(t)\,dt)$.
\end{lemma}
\begin{proof}
Contour shifting in the strip bound of
Lemma~\ref{d:lem:theta-tails} gives
$|\Xi(t)|\le C_a e^{-a|t|}$ for every $0<a<\pi/4$.  Consequently
$d\nu=w|R|^2dt$ has a positive exponential moment.  If an even
$f\in L^2(\nu)$ is orthogonal to all $t^{2j}$, Cauchy--Schwarz gives
a positive exponential moment for the finite complex measure
$f\,d\nu$.  Its Fourier transform is holomorphic on a strip.  All
derivatives at zero vanish, by orthogonality for the even derivatives
and parity for the odd ones.  Analytic continuation and uniqueness of
the Fourier transform of a finite measure imply $f=0$.  Hence the
polynomials in $t^2$ are dense in $L^2_{\mathrm{even}}(\nu)$.
The nonzero entire function $\Xi$ has a real zero set of measure zero.
For $\phi\in L^2_{\mathrm{even}}(w\,dt)$, apply this density to
$\phi/R$, defined off that null set.  Multiplication by $R$ gives
the assertion, without any lower bound on $|\Xi|$.
\end{proof}

\begin{theorem}[Every finite prime head has both infinite indices]
\label{vfourA:thm:indices}
For every integer $M\ge1$, the restrictions of $q_M$ and $p_{>M}$ to
$\Rcal$ have infinite positive and infinite negative index.  Every
witness is an actual full-domain radical of $q_0$.
\end{theorem}
\begin{proof}
Plancherel's identity for the finite head gives, for $u\in\Rcal$,
\begin{equation}\label{vfourA:eq:head-multiplier}
 q_M[u]=\frac1{2\pi}\int_\R g_M(t)|\mathcal Fu(t)|^2\,dt,
 \qquad
 g_M(t)=d_\Gamma(t)-\cg
        -2\sum_{2\le n\le M}\frac{\Lambda(n)}{\sqrt n}\cos(t\log n).
\end{equation}
The continuous even multiplier is negative near zero because
\[
 g_M(0)=-\cg-2\sum_{2\le n\le M}\Lambda(n)/\sqrt n<0.
\]
It tends to $+\infty$ as $|t|\to\infty$, and $|g_M|\le C_Mw$.
Choose any prescribed finite number of orthonormal even smooth
functions supported in a bounded band on which $g_M$ is strictly
negative.  Their multiplier Gram matrix is negative definite.
Lemma~\ref{vfourA:lem:density} approximates each function by a transform
in~\eqref{vfourA:eq:polynomial-radical}.  The bound $|g_M|\le C_Mw$
gives convergence of the whole finite multiplier Gram matrix, while
ordinary $L^2$ convergence preserves independence.  The resulting
actual radical span is negative definite.  The same construction in
a bounded strictly positive band proves the positive assertion.
Finally, the full-domain radical identity and
\eqref{vfourA:eq:head-complete} give
$q_M(u,v)=p_{>M}(u,v)$ for $u,v\in\Rcal$.
\end{proof}

Both infinite indices persist under any finite family of linear
constraints or any finite-rank Hermitian correction: intersect a
definite space of dimension $r+k$ with a subspace of codimension at
most $k$.  Nevertheless $q_0$ vanishes identically on these witnesses.
Their negative head values are therefore cutoff artifacts, not
negative values of the complete form.

\subsection{The complete tail and compactness on the known radical}
\label{vfourA:subsec:compactness}

\begin{lemma}[Complete quantitative tail bound]\label{vfourA:lem:tail}
The raw tail has a bounded self-adjoint representative $D_{>M}$ on
$\Hh$, and
\begin{equation}\label{vfourA:eq:tail-bound}
 \norm{D_{>M}}\le\tau_M
 :=80\sum_{n>M}\Lambda(n)n^{3/2}e^{-\pi n}
 \le160(M+1)^2e^{-\pi(M+1)}.
\end{equation}
Let $D=CD_{>M}C|_{\Hh_c}$.  Then $\norm D\le\tau_M$, and the centered
finite-head representative is
\begin{equation}\label{vfourA:eq:head-operator}
 L_M=L_c+D,\qquad \dom L_M=\dom L_c.
\end{equation}
\end{lemma}
\begin{proof}
The complete theta series gives
$W(x)\le40\exp(4|x|-\pi e^{2|x|})$.
Indeed its coefficient is at most
$4\pi^2\sum_{j\ge1}j^4e^{-\pi(j^2-1)}<40$:
for $j\ge2$ the successive ratio is less than $1/1000$, and
$e^3>20$, $\pi<22/7$ give the upper bound
$4(22/7)^2(1+2/999)<40$.
For $s=\log n$, set $a=|x|$, $b=|x+s|$, $z=a+b\ge s$.
Convexity gives $e^{2a}+e^{2b}\ge2e^z$, whence
\[
 \sup_x\sqrt{W(x)W(x+s)}
 \le40\sup_{z\ge s}e^{2z-\pi e^z}
 =40n^2e^{-\pi n}.
\]
Under the isometry $u\mapsto u/\sqrt W$, each oriented weighted
shift has norm bounded by this supremum.  Pairing both orientations
and summing every $n>M$ proves the first inequality in
\eqref{vfourA:eq:tail-bound}, with absolute operator-norm convergence.
The second follows from $\Lambda(n)\le\log n\le\sqrt n$ and
$\sum_{k\ge0}(k+1)^2e^{-\pi k}<2$.
Centering is performed only in $\Hh$; the raw shifts need not
preserve $\Hh_c$.  Bounded form perturbation proves
\eqref{vfourA:eq:head-operator}.
\end{proof}

Although $\tau_M\to0$, Theorem~\ref{vfourA:thm:indices} retains both
infinite indices.  For example, an ordinary-normalized radical with
$q_M[u]\le-\cg/4$ must satisfy
$\norm u_{\Hh}^2\ge\cg/(4\tau_M)$.  Small operator norm does not
determine sign, and the two normalizations must not be interchanged.

\begin{lemma}[Compact known-radical rows]\label{vfourA:lem:compact-rows}
The maps $DP_\Nn$ and $P_\Nn D$ are compact on $\Hh_c$.
For every sufficiently large $s$, $D(L_c+sI)^{-1}$ is compact.
\end{lemma}
\begin{proof}
The compact embedding $X\hookrightarrow L^2_{\mathrm{even}}(dx)$
is Proposition~\ref{a:prop:compact}.  On $\Nn\subset\ker L_c$,
\[
 \tg[n]=\ip{D_bn}{n}_{\Hh}\le d\norm n_{\Hh}^2,
\]
so $\Nn\hookrightarrow X$ is bounded and
$\Nn\hookrightarrow L^2(dx)$ is compact.  For a fixed shift,
$u\mapsto W(\cdot)u(\cdot+s)$ is bounded from ordinary $L^2$ into
the ambient weighted space, since its squared weighted norm is at
most $\norm W_\infty\norm u_2^2$.  Pair the two orientations to
preserve evenness and then apply $C$ in the weighted space.
Each finite collection of tail rows composed with $P_\Nn$ is
compact.  Lemma~\ref{vfourA:lem:tail} makes $DP_\Nn$ their operator-norm
limit, and adjunction gives compactness of $P_\Nn D$.
Likewise, the shifted-form estimate maps
$(L_c+sI)^{-1}:\Hh_c\to X_c$ boundedly for large $s$; the same
finite-row factorization and norm-tail argument prove the last claim.
No assertion that $C$ or the full infinite tail is bounded from
ordinary $L^2$ to $\Hh$ is needed.
\end{proof}

\subsection{Finite-row repair and the finite-head negative spectrum}
\label{vfourA:subsec:repair}

Let $P_m$ be the weighted orthogonal projection onto
$\operatorname{span}\{r_1,\ldots,r_m\}$, and let $P_\Nn$ project onto
$\Nn$.  Define
\begin{equation}\label{vfourA:eq:repairs}
 F_m=P_mD+DP_m-P_mDP_m,\qquad
 F_\infty=P_\Nn D+DP_\Nn-P_\Nn DP_\Nn.
\end{equation}

\begin{theorem}[Compact repair from actual finite-head rows]
\label{vfourA:thm:repair}
The operator $F_\infty$ is compact and self-adjoint, and
$\rank F_m\le2m$.  With $e_{m,M}=\norm{D(P_\Nn-P_m)}$,
\begin{gather}
 e_{m,M}\longrightarrow0,\qquad
 \norm{F_\infty-F_m}\le2e_{m,M},\label{vfourA:eq:repair-error}\\
 (L_M-F_m)P_m=0,\qquad
 \norm{(L_M-F_m)P_\Nn}\le e_{m,M},\label{vfourA:eq:null-row}\\
 L_M-F_\infty=L_c+P_\Kk DP_\Kk,
 \qquad\dom(L_M-F_\infty)=\dom L_c.\label{vfourA:eq:repaired-block}
\end{gather}
Every finite repair still has infinite positive and negative form
index on $\Rcal$.  The compact compression
$P_\Nn DP_\Nn|_\Nn$ has infinitely many positive and infinitely many
negative eigenvalues, with both sequences tending to zero.
\end{theorem}
\begin{proof}
The projections $P_m$ converge strongly to $P_\Nn$; compactness of
$DP_\Nn$ gives $e_{m,M}\to0$.  The identity
\[
 F_\infty-F_m=(P_\Nn-P_m)D(I-P_m)
                 +(I-P_\Nn)D(P_\Nn-P_m)
\]
proves the norm bound.  Lemma~\ref{vfourA:lem:compact-rows} proves
compactness, and the range of $F_m$ is contained in
$\Ran P_m+D\Ran P_m$, proving its rank bound.
Since $L_cP_\Nn=0$,
\[
 (L_M-F_m)P_\Nn=(I-P_m)D(P_\Nn-P_m),
\]
which gives both null-row assertions.  All products are defined
because $\Nn\subset\dom L_c=\dom L_M$.  Bounded block algebra gives
\eqref{vfourA:eq:repaired-block}.  Finite-rank corrections cannot remove
the two infinite indices of Theorem~\ref{vfourA:thm:indices}.
Finally the compact self-adjoint compression has the same form as
$q_M$ on $\Nn$, so the compact spectral theorem and those indices
give its two eigenvalue sequences.
\end{proof}

The correction uses actual finite-head rows rather than a separately
truncated tail.  Regard $R_m=(r_1,\ldots,r_m)$ as a map
$\C^m\to\Hh_c$, and put $G_m=R_m^*R_m$ and
$Y_m=L_MR_m=DR_m$.  Linear independence from
Lemma~\ref{a:lem:radical} gives $G_m>0$, and
\begin{equation}\label{vfourA:eq:row-formula}
 \begin{split}
 F_m={}&R_mG_m^{-1}Y_m^*+Y_mG_m^{-1}R_m^*\\
       &-R_mG_m^{-1}(R_m^*Y_m)G_m^{-1}R_m^*.
 \end{split}
\end{equation}
Thus $Y_m$ can be formed with the complete Gamma row and only the
finite prime-power head.  Norm convergence supplies no effective
rate for $e_{m,M}$ and no algorithm for the projection onto the full
physical kernel.  The limiting repair leaves $L_c|_\Kk$ unsigned.

\begin{theorem}[Discrete negative head spectrum under the compact-pencil input]
\label{vfourA:thm:negative-spectrum}
Assume the exact ordinary-$L^2$ minimal-form realization input of
Propositions~\ref{a:prop:closedpencil} and~\ref{a:prop:compact}:
for every $\delta>0$, $q_0+\delta\norm{\cdot}_{\Hh}^2$, with form
domain $X$, is lower bounded and has a compact-resolvent
self-adjoint representative on $L^2_{\mathrm{even}}(dx)$.
Then, for every finite integer $M\ge1$,
\[
 \sigma_{\ess}(L_M)\cap(-\infty,0)=\varnothing,
\]
and $L_M$ has infinitely many negative discrete eigenvalues of finite
multiplicity accumulating at zero.
\end{theorem}
\begin{proof}
For any $\delta>0$, the spectral range
$\one_{(-\infty,-\delta]}(L_c)\Hh_c$ lies in $X_c$, since $L_c$ is
lower bounded.  On that range,
\[
 q_0[u]+\frac\delta2\norm u_{\Hh}^2
 \le-\frac\delta2\norm u_{\Hh}^2.
\]
The assumed ordinary compact-resolvent realization has finite negative
index.  The injective inclusion into ordinary $L^2$ therefore bounds
the dimension of this weighted spectral range by that finite index.
This excludes negative essential spectrum of $L_c$; it uses the
ordinary realization to count form signs, not as a unitary model.
By Lemma~\ref{vfourA:lem:compact-rows} and the resolvent identity,
\[
 (L_M+sI)^{-1}-(L_c+sI)^{-1}
 =-(L_M+sI)^{-1}D(L_c+sI)^{-1}
\]
is compact for sufficiently large $s$.  Invariance of essential
spectrum under compact resolvent differences gives the same
essential spectrum for $L_M$ and $L_c$.
The infinite negative index in Theorem~\ref{vfourA:thm:indices}
forces an infinite-dimensional negative spectral subspace of $L_M$.
Lower semiboundedness and the absence of negative essential spectrum
leave isolated finite-multiplicity negative eigenvalues whose only
possible accumulation point is zero.
\end{proof}

If additionally $L_c\ge0$, every negative head eigenvalue lies in
$[-\tau_M,0)$.  Without that extra hypothesis no such enclosure for
\emph{all} head eigenvalues is claimed.  The exact truncation-artifact
statement concerns the radical witnesses; individual negative
eigenvectors of $L_M$ need not belong to $\Nn$.

\subsection{Projection onto a fixed positive spectral range}
\label{vfourA:subsec:projection}

Fix $A>0$ and set
\begin{equation}\label{vfourA:eq:positive-range}
 Q_A=\one_{[A,\infty)}(L_c),\qquad
 T_A=Q_ADQ_A|_{\Ran Q_A}.
\end{equation}
The full Gamma symbol satisfies, for $t\ge1$,
\begin{equation}\label{vfourA:eq:modulation}
 \begin{gathered}
 d_\Gamma(t)=\log t+O(1),\quad |d_\Gamma'(t)|\le C/t,
 \\
 \norm{(D_\Gamma-d_\Gamma(t))(a\cos(t\,\cdot))}_2
 \le C\frac{\log(2+t)}t\norm{a'}_2\quad(t\ge2)
 \end{gathered}
\end{equation}
for every ordinary $H^1$ envelope $a$.
To verify the derivative estimate directly, expand
$d_\Gamma(t)=2\sum_{j\ge0}t^2/[b_j(b_j^2+t^2)]$,
$b_j=2j+1/2$, differentiate, and split the sum at $b_j=t$.
For the modulation estimate, split the Fourier offset $\xi$ at
$|\xi|=t/2$.  The derivative bound in the inner region and the
global logarithmic bound outside give
\[
 |d_\Gamma(t+\xi)-d_\Gamma(t)|
 \le C\log(2+t)|\xi|/t.
\]
Apply Plancherel to each carrier orientation and use evenness of
$d_\Gamma$ for the negative carrier.

\begin{lemma}[Fixed-threshold projection estimate]
\label{vfourA:lem:projection}
For a fixed real even $E\in C_c^\infty(\R)$ let
$u_t=\sqrt W E\cos(t\,\cdot)$.  Then
\begin{equation}\label{vfourA:eq:projection-estimate}
 \norm{Q_ACu_t-u_t}_{\Hh}=O_{E,A}(1/\log t)
 \qquad(t\to\infty).
\end{equation}
Moreover $Q_ACu_t\in\dom L_c\cap\Ran L_c$.
\end{lemma}
\begin{proof}
Set $h_t=(E/\sqrt W)\cos(t\,\cdot)/d_\Gamma(t)$.
It is an actual compact smooth operator-domain vector and
$\norm{h_t}_{\Hh}\le\norm{E/W}_2/d_\Gamma(t)$.
The modulation estimate and $T_\Gamma h_t=WD_\Gamma h_t$ imply
\[
 \norm{T_\Gamma h_t-u_t}_{\Hh}=O_E(1/t),\qquad
 \norm{Lh_t-u_t}_{\Hh}
 \le O_E(1/t)+d\norm{E/W}_2/d_\Gamma(t).
\]
Since $C$ reduces $L$ and $L_c\ge-dI$, the complementary spectral
multiplier $L_c(I-Q_A)$ has norm at most $K_A=\max(d,A)$.
Consequently
\[
 \norm{(I-Q_A)Cu_t}_{\Hh}
 \le O_E(1/t)+(d+K_A)\norm{E/W}_2/d_\Gamma(t).
\]
Integration by parts in the compact smooth moment gives
$\norm{Cu_t-u_t}_{\Hh}=O_E(1/t)$, proving
\eqref{vfourA:eq:projection-estimate}.  Both reducing projections
preserve the operator domain.  Finally the inverse of $L_c$ on
$\Ran Q_A$ is bounded by $1/A$, so every vector in that range
belongs to $\Ran L_c$.
\end{proof}

The envelope and $A$ are fixed during this limit.  No estimate uniform
over shrinking supports is used.

\subsection{Two essential signs of the complete tail}
\label{vfourA:subsec:essential-signs}

Choose any actual prime power $n_0>M$ and write
\begin{equation}\label{vfourA:eq:packet-constant}
 x_0=\tfrac12\log n_0,\qquad
 a_0=\frac{\Lambda(n_0)}{\sqrt{n_0}},\qquad
 A_0=a_0W(x_0)>0.
\end{equation}

\begin{theorem}[Both essential tail signs above every fixed positive threshold]
\label{vfourA:thm:essential-signs}
For every integer $M\ge1$ and every fixed $A>0$,
\begin{equation}\label{vfourA:eq:essential-edges}
 \sup\sigma_{\ess}(T_A)\ge A_0,\qquad
 \inf\sigma_{\ess}(T_A)\le-A_0.
\end{equation}
For every $0<b<A_0$, the spectral subspaces of $T_A$ on
$[b,\infty)$ and $(-\infty,-b]$ are both infinite dimensional.
For each finite $r$, there are $r$-dimensional subspaces of
$\dom L_c\cap\Ran Q_A$ on which, respectively,
$p_{>M}\ge b\norm{\cdot}_{\Hh}^2$ and
$-p_{>M}\ge b\norm{\cdot}_{\Hh}^2$.
Every vector in these spaces has
$q_0[u]\ge A\norm u_{\Hh}^2$ and belongs to $\Ran L_c$.
\end{theorem}
\begin{proof}
Choose real even $\eta\in C_c^\infty((-1,1))$ with $\norm\eta_2=1$,
and set
\[
 \eta_\epsilon(y)=\epsilon^{-1/2}\eta(y/\epsilon),\qquad
 E_\epsilon(x)=\eta_\epsilon(x-x_0)+\eta_\epsilon(x+x_0),\qquad
 u_t=\sqrt W E_\epsilon\cos(t\,\cdot).
\]
Fix a sufficiently small $\epsilon>0$ so the supports are disjoint,
$2\epsilon<\log2$, and the only integer logarithm in
$(\log n_0-2\epsilon,\log n_0+2\epsilon)$ is $\log n_0$.
No positive integer-logarithmic shift correlates a packet with
itself, and only $\log n_0$ correlates the two packets.  Define
\[
 g_\epsilon(y)=\eta_\epsilon(y)^2
           \sqrt{W(x_0+y)W(x_0-y)},\qquad
 A_\epsilon=a_0\int_\R g_\epsilon(y)\,dy.
\]
This envelope is real, even, nonnegative, and smooth, and
$A_\epsilon\to A_0$ as $\epsilon\downarrow0$.
Using both orientations of the selected shift gives the exact
\emph{full-tail} identity
\begin{equation}\label{vfourA:eq:exact-packet}
 p_{>M}[u_t]
 =a_0\left(\cos(t\log n_0)\int g_\epsilon(y)\,dy
       +\int g_\epsilon(y)\cos(2ty)\,dy\right).
\end{equation}
Every other tail term vanishes by support before projection; none
is discarded after projection.  Also
\begin{equation}\label{vfourA:eq:packet-weak}
 \norm{u_t}_{\Hh}^2
 =\int E_\epsilon(x)^2\cos^2(tx)\,dx\longrightarrow1,
 \qquad u_t\rightharpoonup0\quad\hbox{in }\Hh.
\end{equation}
The weak limit is the Riemann--Lebesgue lemma for the $L^1$ product
of the fixed compact envelope with a test vector under the isometry
$u\mapsto u/\sqrt W$.

Choose $t_j\to\infty$ along the phase progressions
$t_j\log n_0\equiv0$ or $\pi\pmod{2\pi}$, respectively.
Equation~\eqref{vfourA:eq:exact-packet} tends to $+A_\epsilon$ or
$-A_\epsilon$.  Lemma~\ref{vfourA:lem:projection} gives
$v_j=Q_ACu_{t_j}=u_{t_j}+o_{\Hh}(1)$.  The complete norm-tail bound
of Lemma~\ref{vfourA:lem:tail} preserves the quadratic limits, so
\[
 \norm{v_j}_{\Hh}\longrightarrow1,\qquad
 v_j\rightharpoonup0,\qquad
 \frac{\ip{T_Av_j}{v_j}_{\Hh}}{\norm{v_j}_{\Hh}^2}
 \longrightarrow\pm A_\epsilon.
\]
For a bounded self-adjoint operator, a normalized weakly null
sequence has quadratic limsup at most the top of the essential
spectrum and liminf at least its bottom.  Indeed, the spectral
projection beyond either essential edge by any fixed positive
distance has finite rank and its norm on such a sequence tends
to zero.  Thus the essential edges bound $\pm A_\epsilon$.
Letting $\epsilon\downarrow0$ \emph{after each fixed-envelope limit}
proves~\eqref{vfourA:eq:essential-edges} and both infinite spectral
dimensions.

To obtain actual operator-domain definite spaces, choose
$b<b'<A_0$ and $r$ independent vectors in the spectral subspace of
$T_A$ on $[b',\infty)$, or on $(-\infty,-b']$.
The set $\dom L_c\cap\Ran Q_A$ is dense in $\Ran Q_A$ because
$Q_A$ preserves $\dom L_c$.  Approximate the finite basis from
that intersection.  Continuity of both finite Gram matrices
preserves independence and the strict $\pm p_{>M}\ge
b\norm{\cdot}_{\Hh}^2$ inequality.  No invariance of $\dom L_c$
under the spectral projections of $T_A$ is needed.
The complete-energy inequality and range assertion follow from
the definition of $Q_A$ and Lemma~\ref{vfourA:lem:projection}.
\end{proof}

\subsection{Kernel removal, compact corrections, and the low-energy boundary}
\label{vfourA:subsec:scope}

Because $\Nn\subseteq\ker L_c$, $\Ran Q_A$ lies in both $\Kk$ and
$\Hh_c\ominus\ker L_c$.  The same weakly null sequences show both
essential signs for the tail compression to either of these larger
spaces.  Any compact self-adjoint correction has quadratic values
tending to zero on these sequences.  Hence neither removal of the
known radical nor removal of the entire physical kernel, and no
compact correction on the displayed spaces, makes the compressed
tail one-sided or compact.  In particular,
$P_\Kk F_\infty P_\Kk=0$, so the known-radical repair leaves the
tail quadratic form on $\Kk$ unchanged.

One may take $n_0=2^k>M$.  The first positive theta summand and
$e^{x/2}/(2\cosh(x/2))\ge1/2$ for $x\ge0$ give
\[
 W((\log n_0)/2)\ge\pi^2n_0^2e^{-\pi n_0},\qquad
 A_0\ge\pi^2(\log2)n_0^{3/2}e^{-\pi n_0}.
\]
Here $4\pi^2n^2-6\pi n\ge2\pi^2n^2$ for $n\ge2$.
Thus the essential norm after any compact repair is at least the
displayed positive constant.  This is a bound for the chosen prime
power, not a uniform positive fraction of $\tau_M$ as $M$ varies.

Theorem~\ref{vfourA:thm:essential-signs} holds at every prescribed
fixed threshold.  In particular its witnesses can be constructed
in $\Ran Q_H$ for arbitrarily large $H>0$, where
\[
 q_0[u]\ge H\norm u_{\Hh}^2,\qquad
 \frac{|p_{>M}[u]|}{q_0[u]}\le\frac{\tau_M}{H}
 \quad(u\ne0).
\]
This is a high-energy \emph{absolute-sign} obstruction, compatible
with arbitrarily small relative tail cost on those witnesses.  It
does not determine a tail-to-energy sign or bound in a low-energy
collar near zero.  In particular, at $H>\tau_M$ these same vectors
even have $q_M[u]>0$, regardless of the sign of their tail value.
The radical result concerns exact zero-energy directions instead.
Neither result identifies $\Nn$ with the full kernel, resolves the
complete sign on $\Kk$, or proves or disproves the Riemann hypothesis.
